\section{Conclusion}\label{sec:conclusion}
%This paper examines vector quantization (VQ) from a distributional perspective, introducing three key evaluation criteria. Empirical results demonstrate that optimal VQ results are achieved when the distributions of continuous feature vectors and code vectors are identical. Our theoretical analysis corroborates this finding, emphasizing the crucial role of distributional alignment in effective VQ. Based on these insights, we propose using the quadratic Wasserstein distance to achieve alignment, leveraging its computational efficiency under a Gaussian hypothesis. This approach achieves near-full codebook utilization while significantly reducing quantization error. Our method successfully addresses both training instability and codebook collapse, leading to improved downstream image reconstruction performance. A limitation of this work, 

%This paper examines vector quantization (VQ) from a distributional perspective, introducing three key evaluation criteria. Empirical results show that optimal VQ performance is achieved when the distributions of continuous feature vectors and code vectors are identical. Our theoretical analysis corroborates this finding, highlighting the critical role of distributional alignment in effective VQ. Motivated by these insights, we propose leveraging the quadratic Wasserstein distance to achieve alignment, taking advantage of its computational efficiency under a Gaussian assumption. This approach enables near-complete codebook utilization while substantially reducing quantization error. Our method successfully mitigates both training instability and codebook collapse, leading to improved downstream image reconstruction performance. One limitation of this work 

We identify a fundamental distributional mismatch between feature and code vectors as a common cause of training instability and codebook collapse in vector quantization. To address this issue, we propose a general distributional matching framework that formalizes desirable VQ behavior through principled criteria. Within this framework, we show that aligning feature and code distributions provides a unified mechanism for improving training stability, codebook utilization, and representation efficiency. We instantiate the framework using a Wasserstein-based objective with an efficient closed-form solution under a mild Gaussian approximation, and further demonstrate that a nonparametric alternative based on maximum mean discrepancy achieves comparable performance. More broadly, this work offers a unified perspective on vector quantization and opens new avenues for designing stable and efficient discrete representation learning methods.

\section{Limitation and Discussion}
One limitation of this work is that we do not directly validate the effectiveness of the proposed distributional matching framework through downstream generative results, due to limited computational resources. Recent studies have shown that improved reconstruction performance does not necessarily translate into better generative performance. For instance, increasing the codebook size or token length in discrete visual tokenizers~\cite{yu2024language}, or increasing the latent resolution or dimensionality in continuous tokenizers, can significantly improve reconstruction quality while simultaneously degrading generative performance~\cite{Yao2025ReconstructionVG}. 

In these cases, larger codebooks or longer token sequences in discrete tokenizers substantially increase the complexity of training autoregressive models, due to longer sequences and larger vocabularies. Similarly, higher latent resolution or dimensionality in continuous tokenizers makes the training of diffusion models more challenging, leading to reduced generative quality. 

In contrast, distributional matching improves tokenizer quality without introducing additional burdens to downstream generative models. For example, compared with VAR~\cite{Tian2024VisualAM}, Wasserstein VAR achieves improved reconstruction without increasing the codebook size or token length. As a result, the improved reconstruction performance induced by distributional matching is more likely to translate into improved generative performance.